\documentclass[12pt]{article}
\usepackage[utf8]{inputenc}
\usepackage[margin=0.75in]{geometry}
\usepackage{amsmath}
\usepackage{amssymb}
\usepackage{amsthm}
\usepackage{graphicx} % Required for inserting images
\usepackage{array}
\usepackage{tikz-cd}


\title{MAT1100 Homework 1}
\author{Aaron Avram}
\date{September 20 2026}

\begin{document}
\newtheorem{definition}{Definition}
\newtheorem{theorem}{Theorem}
\newtheorem{claim}{Claim}
\newtheorem{corollary}{Corollary}

\maketitle

\begin{align*}
    \textbf{Question 6}
\end{align*}

\begin{enumerate}
    \item[(a)]
    \begin{proof}
    Since $G/Z(G)$ is cyclic it is generated by a single coset. Let $z$ be a representative of this coset.
    Then all cosets of $Z(G)$ in $G$ are of the form $z^aZ(G)$ for some $a \in \mathbb Z$.

    Now consider $x, y \in G$ they each lie in some coset of $Z(G)$, therefore there exist $a, b \in \mathbb Z$
    such that $x \in z^aZ(G)$ and $y \in z^bZ(G)$. Therefore there exist $g, h \in Z(G)$ such that
    $x = z^ag$ and $y = z^bh$. Now compute
    \begin{align*}
        xy &= z^agz^bh \\
        &\text{since $g \in Z(G)$:} \\
        &= z^az^bgh \\
        &\text{since $h \in Z(G)$:} \\
        &= z^{a+b}hg \\
        &= z^bz^ahg \\
        &\text{since $h \in Z(G)$:} \\
        &= z^bhz^ag \\
        &= yx
    \end{align*}
    hence $G$ is abelian.
    \end{proof}

    \item[(b)]
    \begin{proof}
        Let $p$ be prime and $G$ a group of order $p^2$. Let $G \curvearrowright G$
        by conjugation, then let $c(x_1), \ldots, c(x_k)$ be the distinct non-trivial conjugacy classes
        of $G$, hence $|c(x_i)| > 1$. By orbit-stabilizer:
        \[ |c(x_i)| = |G|/|G_{x_i}| = p^2/|G_{x_i} \]
        since $|c(x_i)| > 1$, $|G_{x_i}| < p^2$ which implies that $|G_{x_i}| \in \{ 1, p \}$.
        Now if $|G_{x_i}| = 1$ then $|c(x_i)| = |G|$ which implies $e \in c(x_i)$.
        Since conjugacy classes are disjoint or equal this forces $c(x_i) = c(e)$
        which means $|c(x_i)| = 1$ contradicting $|c(x_i)| = |G| > 1 $.

        Thus $|G_{x_i}| = p$ and therefore $|c(x_i)| = p$. Now writing the class equation of $G$
        \[ |G| = |Z(G)| + \sum_{i = 1}^k |c(x_i)| \]
        then mod p, $|G| = 0$ and $|c(x_i)| = 0$ therefore this implies
        \[ 0 = |Z(G)| \mod p \]
        now together with Lagrange's theorem this implies $|Z(G)| \in \{p, p^2 \}$.

        Case 1: $|Z(G)| = p^2$ then $Z(G) = G$ and then clearly $G$ is abelian
        
        Case 2: $|Z(G)| = p$ then $|G/Z(G)| = p$. Picking any non identity
        element $g \in G/Z(G)$ (which exists as $p > 1$) we must have $|g| = p$ by Lagrange's theorem
        and hence $G/Z(G) = \langle g \rangle$ meaning it is cyclic. Then by part (a),
        $G$ is abelian.
    \end{proof}

    \item[(c)]
    \begin{proof}
        Consider $G = D_8$. It has a cyclic normal subgroup $R = \langle r \rangle$,
        this is clear as $D_8 = \langle r, s \rangle$ and for any $r^a$ we have
        \[ sr^as^{-1} = r^{-a}ss^{-1} = r^{-a} \in R \]
        then $|D_8/R| = 2$ which implies it is cyclic. However $D_8$ is not abelian
        as
        \[ sr = r^{-1}s \]
        this is not a contradiction as $Z(G) = \langle e, r^2 \rangle$ which is a proper subgroup of $R$.
        And morever $D_8/Z(G)$ is isomorphic to the Klein 4-group which is not cyclic.
    \end{proof}
\end{enumerate}

\end{document}